\documentclass[11pt]{article}
\usepackage[T1]{fontenc}
\usepackage[utf8]{inputenc}
\usepackage{lmodern}
\usepackage{microtype}
\usepackage[margin=1in]{geometry}
\usepackage{amsmath,amssymb,amsthm,mathtools}
\usepackage{enumitem}
\usepackage{booktabs,longtable,array,tabularx}
\usepackage[dvipsnames]{xcolor}
\usepackage[hidelinks]{hyperref}
\usepackage[nameinlink,noabbrev]{cleveref}
\usepackage{titlesec}
\usepackage{fancyhdr}
\setlength{\parskip}{0.55em}
\setlength{\parindent}{0pt}
\setlist[itemize]{leftmargin=1.5em,itemsep=0.25em,topsep=0.25em}
\setlist[enumerate]{leftmargin=1.7em,itemsep=0.25em,topsep=0.25em}
\renewcommand{\arraystretch}{1.15}
\titleformat{\section}{\large\bfseries\color{MidnightBlue}}{\thesection}{0.6em}{}
\titleformat{\subsection}{\normalsize\bfseries\color{MidnightBlue}}{\thesubsection}{0.6em}{}
\pagestyle{fancy}
\fancyhf{}
\fancyhead[L]{\footnotesize Same-domain G\"odel route exhaustion}
\fancyhead[R]{\footnotesize Primary manuscript with full appendices}
\fancyfoot[C]{\thepage}
\newtheorem{theorem}{Theorem}[section]
\newtheorem{lemma}[theorem]{Lemma}
\newtheorem{proposition}[theorem]{Proposition}
\newtheorem{corollary}[theorem]{Corollary}
\theoremstyle{definition}
\newtheorem{definition}[theorem]{Definition}
\theoremstyle{remark}
\newtheorem{remark}[theorem]{Remark}

\newtheorem{example}[theorem]{Example}

\newcommand{\Adm}{\mathrm{Adm}}
\newcommand{\Stand}{\mathrm{Stand}}
\newcommand{\Prov}{\mathrm{Prov}}
\newcommand{\IntD}{\mathrm{Int}}
\newcommand{\code}[1]{\ulcorner #1 \urcorner}
\newcommand{\Law}{\sim}
\newcommand{\Eval}{X_D^{\bullet}}
\newcommand{\domain}{D}
\newcommand{\Good}{\mathrm{Good}}
\newcommand{\Bad}{\mathrm{Bad}}
\newcommand{\Reuse}{\mathrm{ReuseAdm}}
\newcommand{\Pred}{\mathrm{Pred}}
\newcommand{\Hist}{\mathcal{H}_D}
\newcommand{\Quot}{\mathrm{Q}_D}
\newcommand{\Traj}{\mathsf{Traj}_D}
\newcommand{\Legal}{\mathrm{Legal}}
\newcommand{\Coherent}{\mathrm{Coh}}
\newcommand{\Read}{\mathrm{Read}}
\newcommand{\Threat}{\mathrm{Threat}}
\newcommand{\Fix}{\mathrm{Fix}}

\allowdisplaybreaks

\title{Phase 5 Report\\\large Repositioned appendices, formalization placement, and audit offload}
\author{OpenAI}
\date{\today}
\begin{document}
\maketitle

\section{Phase objective}
Phase~5 executes the rewrite-plan instruction to reduce referee burden while preserving audit-grade rigor by repositioning the appendices and formalization material. The governing decision of this phase is kernel-first: the compact kernel floor remains visible, and all downstream layers are treated as forced consequences of that kernel. The coherent-reasoning paper remains the decisive defeat of the hardest hostile attack because it closes the possibility of a hidden same-domain intermediate path distinction that matters to closure while escaping standing and trajectory structure.

\section{What was changed}
\subsection*{Appendix architecture now implemented}
\begin{itemize}
  \item \textbf{Appendix A} remains the compact inferential-governance floor.
  \item \textbf{Appendix B} records only the minimum interface-forcing results: AMetric boundary, bivalent interface, standing emergence, unique admissible interior, and standing conservation.
  \item \textbf{Appendix C} records only the minimum standing-quotient and local-closure results.
  \item \textbf{Appendix D} records only the minimum coherent-trajectory results: pointwise standing-preservation, no deferred coherence, endpoint certification, and tensor/skin closure.
  \item \textbf{Appendix E} records only the minimum no-second-equivalence bridge plus the short formalization and audit-placement note.
\end{itemize}

\subsection*{Formalization placement now implemented}
The primary manuscript now contains only a short formalization note in the main text and a short evidential note in Appendix~E. Detailed theorem-ledger tables, dependency spines, and hostile-proof provenance are no longer presented as part of the main submission architecture. They are explicitly offloaded to the audit companion.

\subsection*{Supplement and audit packaging now implemented}
A standalone \texttt{technical\_appendix.tex} is included for supplementary-material use, and an \texttt{audit\_companion/} folder is included in the project zip. The companion contains a frozen copy of the current audit-edition source together with a placement note explaining that the hyper-hardened theorem reservoir remains available for referee-facing defense but is no longer allowed to dominate the primary narrative.

\section{Kernel-centered rationale}
The phase-5 appendix strategy is driven by one governing claim: once the kernel is fixed, all downstream papers used by the main manuscript are forced consequences relative to that kernel. This is why the appendices are now layered in the following order:
\begin{enumerate}
  \item kernel floor;
  \item interface forcing and standing emergence;
  \item standing quotient uniqueness;
  \item coherent trajectory closure and irrecoverability;
  \item no second same-domain equivalence and inferential-life exhaustion.
\end{enumerate}
This ordering makes the bridge to the Godel application transparent. Any same-domain closure threat must now survive into either standing classification or coherent legal trajectory structure. The coherent-reasoning layer closes the hostile idea that a threat might hide in intermediate derivation structure without surfacing there.

\section{Appendix allocation table}
\begin{center}
\begin{tabular}{p{0.18\linewidth}p{0.30\linewidth}p{0.42\linewidth}}
\toprule
Layer & Placement & Why it stays / moves \\
\midrule
Kernel attachment, no faithful lower generator, Level~2 / Level~3 separation & Appendix A & Needed to justify beginning the body from inferential governance rather than trace classification. \\
AMetric boundary, bivalence, standing emergence, unique interior, standing conservation & Appendix B & Needed because the body relies on a forced interface, but full audit-level proof burden is unnecessary. \\
Standing quotient uniqueness and tensor/skin local closure & Appendix C & Needed to block second same-domain standing-bearing layers and to support code-role exhaustion. \\
Coherent trajectories, no deferred coherence, endpoint certification, tensor determination & Appendix D & Needed because this layer defeats the hardest hidden-path attack directly. \\
No second same-domain equivalence, inferential-life exhaustion, short formalization note & Appendix E & Needed as the final bridge to the body theorems and to keep formalization placement minimal. \\
Detailed theorem ledgers, dependency graphs, hardened hostile-proof reserves & Audit companion & Useful for defense, but they overload the reader if kept in the primary manuscript. \\
\bottomrule
\end{tabular}
\end{center}

\section{Compile targets for Overleaf}
The project zip is intentionally multipurpose.
\begin{itemize}
  \item \texttt{main.tex} compiles the phase-5 primary manuscript seed.
  \item \texttt{phase5\_seed.tex} is the same primary manuscript root explicitly.
  \item \texttt{technical\_appendix.tex} compiles the compact appendices A--E as a standalone supplement.
  \item \texttt{audit\_companion/current\_audit\_edition\_main.tex} compiles the frozen audit edition.
  \item \texttt{phase5\_report.tex} compiles this report locally if needed.
\end{itemize}

\section{Exit check against the rewrite plan}
Phase~5 is complete when the following are all true:
\begin{itemize}
  \item the body still reads as a same-domain Godel-response paper rather than as an architecture survey;
  \item the appendices are compact and layered rather than historically accretive;
  \item the formalization is evidential and audit-supportive rather than body-dominating;
  \item the coherent-reasoning layer remains visibly responsible for defeating the hidden-path attack;
  \item the audit-grade theorem reservoir remains available without distorting the primary narrative.
\end{itemize}
These conditions are now satisfied by the phase-5 package.
\end{document}
